\subsection{Vision Transformer}
\label{sec:vision_transformer}

\begin{figure}[htbp]
\centering
\includegraphics[width=\columnwidth]{vit.png}
\caption{Architecture of the ViT Backbone and Patch Projection.}
\label{fig:3}
\end{figure}

The visual representation extraction architecture of this study adopts ViT-Base, as illustrated in Figure~\ref{fig:3} \cite{ref26}. An input facial image of $224 \times 224$ pixels with 3 color channels is partitioned into 196 non-overlapping spatial patches $\mathbf{x}_p^i$ of size $16 \times 16$ pixels \cite{ref27}. Each patch is linearly projected via matrix $\mathbf{E}$ into a feature representation space of dimension $D = 768$, and subsequently concatenated with a class token $\mathbf{x}_{\text{class}}$ and positional embedding $\mathbf{E}_{\text{pos}}$ according to \eqref{eq:1} \cite{ref27}. The initial token sequence vector $\mathbf{z}_0$ is fed into a stack of 12 identical transformer encoder layers. Each encoder layer $\ell$ processes the input $\mathbf{z}_{\ell-1}$ into an intermediate representation $\mathbf{z}'_\ell$ via the Multi-Head Self-Attention (MHSA) mechanism following a Layer Normalization (LN) operation according to \eqref{eq:2}, which then proceeds to an output representation $\mathbf{z}_\ell$ through a two-layer Multi-Layer Perceptron (MLP) block with Gaussian Error Linear Unit (GeLU) activation based on \eqref{eq:3} \cite{ref28}. The self-attention mechanism enables the modeling of global relationships across facial patches without dependence on local receptive fields as in convolutional operations \cite{ref29}.

\begin{equation}
\mathbf{z}_0 = [\mathbf{x}_{\text{class}}; \, \mathbf{x}_p^1\mathbf{E}; \, \dots; \, \mathbf{x}_p^{196}\mathbf{E}] + \mathbf{E}_{\text{pos}}
\label{eq:1}
\end{equation}

\begin{equation}
\mathbf{z}'_\ell = \text{MHSA}(\text{LN}(\mathbf{z}_{\ell-1})) + \mathbf{z}_{\ell-1}
\label{eq:2}
\end{equation}

\begin{equation}
\mathbf{z}_\ell = \text{MLP}(\text{LN}(\mathbf{z}'_\ell)) + \mathbf{z}'_\ell
\label{eq:3}
\end{equation}

\begin{table}[htbp]
\caption{Multi-Domain Feature Fusion and Ablation Configurations.}
\label{tab:II}
\centering
\begin{tabular}{clcc}
\toprule
\# & Configuration & Domain Category & Dimension \\
\midrule
1 & Face & Single-Domain & 768 \\
2 & Emotion & Single-Domain & 768 \\
3 & Age & Single-Domain & 768 \\
4 & Emotion $\oplus$ Face & Dual-Domain & 1,536 \\
5 & Face $\oplus$ Age & Dual-Domain & 1,536 \\
6 & Emotion $\oplus$ Age & Dual-Domain & 1,536 \\
7 & Face $\oplus$ Emotion $\oplus$ Age & Tri-Domain (Proposed) & 2,304 \\
\bottomrule
\end{tabular}
\end{table}

To maintain computational efficiency and prevent retraining variability, the three ViT backbone models are frozen as providers of task-associated representations via offline extraction. The ViT-Face model (skutaada/\allowbreak VIT-\allowbreak VGGFace) captures representations associated with facial biometric geometry, ViT-Emotion (dima806/\allowbreak facial\_\allowbreak emotions\_\allowbreak image\_\allowbreak detection) produces representations associated with facial affective expression, and ViT-Age (dima806/\allowbreak facial\_\allowbreak age\_\allowbreak image\_\allowbreak detection) extracts representations associated with facial age estimation. From each domain model, the 768-dimensional visual representation feature vector $\mathbf{f}_{\text{domain}} \in \mathbb{R}^{768}$ is extracted from the $[\text{CLS}]$ token representation at the final encoder layer $L$ ($\mathbf{z}_L^0$) following the LN operation according to \eqref{eq:4}. These three single-domain feature representations are combined into a tri-domain fusion vector $\mathbf{z}_{\text{tri}} \in \mathbb{R}^{2304}$ via a feature concatenation operation ($\oplus$) of the vectors $\mathbf{f}_{\text{face}}$, $\mathbf{f}_{\text{emotion}}$, and $\mathbf{f}_{\text{age}}$ as formulated in \eqref{eq:5}. The systematic exploration encompasses seven feature ablation schemes, comprising three 768-dimensional single-domain configurations, three 1,536-dimensional dual-domain configurations, and one 2,304-dimensional tri-domain configuration, as summarized in Table~\ref{tab:II}.

\begin{equation}
\mathbf{f}_{\text{domain}} = \text{LN}(\mathbf{z}_L^0) \in \mathbb{R}^{768}
\label{eq:4}
\end{equation}

\begin{equation}
\mathbf{z}_{\text{tri}} = \mathbf{f}_{\text{face}} \oplus \mathbf{f}_{\text{emotion}} \oplus \mathbf{f}_{\text{age}} \in \mathbb{R}^{2304}
\label{eq:5}
\end{equation}
